\documentclass[a4paper,twocolumn,11pt,unpublished,nopdfoutputerror]{quantumarticle}
\usepackage[utf8]{inputenc}
\usepackage[T1]{fontenc}
\usepackage[english]{babel}
\usepackage{amsmath,amssymb,amsthm,mathtools}
\usepackage{braket}
\usepackage{graphicx}
\usepackage{booktabs}
\usepackage{siunitx}
\usepackage{microtype}
\usepackage{xcolor}
\usepackage[numbers,sort&compress]{natbib}
\usepackage[ruled,vlined]{algorithm2e}
\usepackage{enumitem}
\usepackage{multirow}
\usepackage{bbm}

% ---------- macros ----------
\DeclareMathOperator{\Tr}{Tr}
\newcommand{\F}{\mathbb{F}}
\newcommand{\R}{\mathbb{R}}
\newcommand{\Pauli}{\mathcal{P}}
\newcommand{\synd}{\mathbf{s}}
\newcommand{\err}{\mathbf{e}}
\newcommand{\ehat}{\hat{\mathbf{e}}}
\newcommand{\llr}{\Lambda}
\newcommand{\Hmat}{\mathbf{H}}
\newcommand{\Gmat}{\mathbf{G}}
\newcommand{\Wmat}{\mathbf{W}}
\newcommand{\bvec}[1]{\mathbf{#1}}
\newcommand{\phat}{\hat{p}}
\newcommand{\concat}{\,\|\,}
\newcommand{\code}[3]{\ensuremath{[\![#1,#2,#3]\!]}}

\usepackage{graphicx}
\graphicspath{{./figures/}}
\DeclareGraphicsExtensions{.pdf,.png,.jpg}

% Placeholder for missing data
\newcommand{\pending}[1]{\textcolor{red}{\textbf{[#1]}}}

% Review comment macros (remove before submission)
\newcommand{\anish}[1]{{\small\color{purple}[Anish] #1}}
\newcommand{\sai}[1]{{\small\color{teal}[Sai] #1}}
\newcommand{\nithin}[1]{\textcolor{blue}{[Nithin: #1]}}

\title{Adaptive GNN-Enhanced Belief Propagation Decoding of Quantum LDPC Codes Under Realistic Noise}

\author{Jothiradithya Sai Konduru}
\affiliation{Paradise Valley High School, Phoenix, AZ, USA}

\author{Anish Chedalla}
\affiliation{Paradise Valley High School, Phoenix, AZ, USA}

\author{Nithin Raveendran}
\affiliation{Department of Electrical and Computer Engineering, University of Arizona, Tucson, AZ, USA}
\email{nithin@arizona.edu}

\thanks{The authors thank the STAR Lab program at the University of Arizona for facilitating this research.}

\begin{document}
\maketitle

% ============================================================
%  ABSTRACT
% ============================================================
\begin{abstract}
Fault-tolerant quantum computing depends crucially on classical decoding algorithms that infer and correct errors from syndrome measurements. Quantum low-density parity-check (QLDPC) codes offer dramatically improved encoding rates over surface codes, yet their short-cycle Tanner graph structure causes standard belief propagation (BP) to fail on the majority of decoding instances. This problem is compounded by noise non-stationarity on real hardware, where biased error rates drift over time and violate fixed-parameter decoder assumptions. We propose a hybrid decoding framework for bivariate bicycle QLDPC codes that augments---rather than replaces---BP with a graph neural network (GNN) operating directly on the code's Tanner graph. The framework introduces four integrated components: (i)~\textsc{TannerGNN}, a GNN with edge-type-aware message passing that produces per-qubit log-likelihood ratio corrections; (ii)~Feature-wise Linear Modulation (FiLM) conditioning that enables a single trained model to adapt to varying noise rates without retraining; (iii)~a fully differentiable neural BP decoder with learnable per-iteration message-scaling weights; and (iv)~an interleaved architecture in which the GNN intervenes mid-decoding using intermediate BP marginals. We evaluate nine decoder configurations across three bivariate bicycle codes ($\code{72}{12}{6}$, $\code{144}{12}{12}$, $\code{288}{12}{18}$) under $Z$-biased noise ($\eta=20$) with static and temporally drifting noise, supported by component ablation, transfer-learning analysis, and McNemar paired significance testing. \pending{DATA: Insert headline LER reduction number once runs complete.}
\end{abstract}


% ============================================================
%  I. INTRODUCTION  (Nithin's text, preserved)
% ============================================================
\section{Introduction}

\subsection{Context and problem}

Quantum computing has attracted significant attention due to demonstrations of computational advantage in controlled experiments \cite{Arute2019Supremacy,Zhong2020PhotonicAdvantage}, alongside its potential impact on applications such as quantum chemistry \cite{Wang2008QuantumAlgorithm}. A central obstacle to practical quantum computation is the fragility of quantum information: quantum states are easily corrupted by interactions with the environment, and these errors accumulate as circuits grow in depth. Quantum error correction (QEC) addresses this challenge by encoding logical information across many physical qubits so that errors can be detected and corrected without directly measuring the encoded state \cite{Shor1995Scheme,Lidar2013Background}. In the stabilizer formalism, error information is extracted via ancilla-assisted stabilizer measurements that report parity constraints; aggregating these binary outcomes produces a syndrome that indicates which checks are violated \cite{Gottesman1997Stabilizer}.

While QEC enables fault tolerance in principle, the overhead of many traditional codes can be prohibitive, motivating interest in quantum low-density parity-check (QLDPC) codes. QLDPC codes leverage sparse parity-check structure to support efficient scaling, and recent developments suggest that QLDPC families can achieve high thresholds with reduced overhead compared to earlier approaches \cite{Bravyi2024Memory,Panteleev2021QLDPC}. However, decoding remains a critical bottleneck. The standard approach for LDPC-style decoding is belief propagation (BP), an iterative message-passing algorithm on the Tanner graph induced by the parity-check matrix \cite{Roffe2020Landscape}. BP is attractive because of its low per-iteration complexity and amenability to hardware-friendly implementations, but its performance depends strongly on assumed error priors and independence structure. These assumptions are often violated by realistic device noise, where error types can be strongly biased (e.g., dephasing-dominated noise) and drift over time due to calibration changes or environmental fluctuations \cite{Tuckett2018Bias}. As a result, BP can perform poorly under prior mismatch even when the syndrome contains sufficient information to recover the correct logical state.

Recent work has shown that learned decoders can improve performance on difficult QLDPC Tanner graphs. Maan and Paler~\cite{Maan2025MLBP} developed Astra, a scalable GNN decoder that replaces BP entirely, achieving higher thresholds than BP-OSD on bivariate bicycle codes up to distance~18. Ninkovi\'{c} et al.~\cite{Ninkovic2024GNN} proposed a GNN decoder exploiting the sparse graph structure of QLDPC codes with linear complexity scaling. Gong et al.~\cite{Gong2024GNN} constructed a fully differentiable iterative decoder that interleaves BP stages with intermediate GNN layers, lowering the error floor on quantum bicycle codes. For repetition codes on IBM hardware, Stein et al.~\cite{Stein2026FiLM} applied Feature-wise Linear Modulation (FiLM) conditioning to enable noise-adaptive decoding, and Blue et al.~\cite{Blue2026MLBB} demonstrated a recurrent transformer that outperforms BP-OSD by a factor of~5 on the $\code{72}{12}{6}$ code at circuit level. In classical coding, Nachmani et al.~\cite{Nachmani2016Learning,Nachmani2018Deep} showed that learnable per-iteration edge weights can compensate for short cycles in BP. However, these approaches either replace BP entirely, address only surface codes, or operate under simplified noise models. No prior work has combined GNN-based LLR correction, FiLM noise adaptivity, differentiable neural BP, and mid-decoding interleaved correction within a single framework for QLDPC codes.

\subsection{Contributions}

We propose a hybrid decoding framework that augments belief propagation with data-driven noise adaptation for bivariate bicycle (BB) CSS-type QLDPC codes. The core idea is to separate noise inference from structured parity-constrained decoding: a graph neural network processes the Tanner graph structure together with observed syndrome information to infer adaptive per-qubit error priors, which are then provided to a BP decoder that performs iterative message passing to produce a correction \cite{Roffe2020Landscape}. This design preserves compatibility with BP while allowing the learned component to compensate for biased and temporally drifting noise regimes that violate fixed-prior assumptions \cite{Tuckett2018Bias}.

Specifically, we make the following contributions:
\begin{enumerate}[leftmargin=*,topsep=2pt,itemsep=2pt]
\item \textsc{TannerGNN}, a graph neural network with edge-type-aware message passing on the CSS Tanner graph, producing per-qubit LLR corrections that substantially improve BP convergence (Sec.~\ref{sec:tannergnn}).

\item Feature-wise Linear Modulation (FiLM) conditioning, applied for the first time to QLDPC decoding, enabling a single model to maintain performance across a range of noise rates without retraining (Sec.~\ref{sec:film}).

\item A fully differentiable neural BP decoder with learnable per-iteration message-scaling weights, trained end-to-end jointly with the GNN through gradient descent (Sec.~\ref{sec:neuralbp}).

\item An interleaved GNN-BP architecture---to our knowledge the first for any quantum code family---in which the GNN intervenes mid-decoding using intermediate BP marginals (Sec.~\ref{sec:interleaved}).

\item Comprehensive evaluation across three BB codes, four noise regimes, and nine decoder configurations, with component ablation, transfer-learning analysis, and McNemar significance testing (Sec.~\ref{sec:results}).
\end{enumerate}

We note that while reinforcement learning has been proposed for adaptive prior tuning in iterative decoders, we find that end-to-end differentiable training through BP provides direct per-qubit gradient information, avoiding the credit assignment and sample efficiency challenges inherent to policy-gradient methods. FiLM conditioning further renders online RL unnecessary by enabling noise adaptation through a single conditioned forward pass.

\subsection{Organization of the paper}
\begin{itemize}
  \item Section~\ref{sec:background}: QEC background, CSS codes, bivariate bicycle construction, BP decoding, short-cycle problem, noise models
  \item Section~\ref{sec:methods}: proposed framework---TannerGNN, FiLM conditioning, neural BP, interleaved architecture, training
  \item Section~\ref{sec:setup}: experimental setup, code parameters, hyperparameters, metrics
  \item Section~\ref{sec:results}: simulation results, threshold curves, drift robustness, ablation, transfer learning
  \item Section~\ref{sec:discussion}: comparison with prior work, limitations, design philosophy
  \item Section~\ref{sec:conclusion}: summary and future directions
\end{itemize}


% ============================================================
%  II. BACKGROUND AND PRELIMINARIES
% ============================================================
\section{Background and preliminaries}\label{sec:background}

\subsection{Stabilizer codes and the CSS construction}

A Calderbank--Shor--Steane (CSS) code~\cite{Calderbank1996Good,Steane1996Error} is defined by a pair of binary parity-check matrices $\Hmat_x \in \F_2^{m_x \times n}$ and $\Hmat_z \in \F_2^{m_z \times n}$ satisfying the commutativity constraint:
\begin{equation}
\Hmat_x \Hmat_z^\top = \mathbf{0} \pmod{2}.
\label{eq:css}
\end{equation}
The rows of $\Hmat_x$ define $X$-type stabilizers that detect $Z$~errors, and vice versa. Given an error $\err = (\err_x, \err_z)$, the syndrome is:
\begin{equation}
\synd_x = \Hmat_x \err_z \pmod{2}, \quad \synd_z = \Hmat_z \err_x \pmod{2}.
\label{eq:syndrome}
\end{equation}
Decoding amounts to finding an estimate $\ehat$ consistent with the observed syndrome and equivalent to $\err$ up to stabilizer equivalence. A code with $n$~physical qubits, $k$~logical qubits, and minimum distance~$d$ is denoted $\code{n}{k}{d}$, with encoding rate $k/n$.

\subsection{Bivariate bicycle codes}

Bivariate bicycle (BB) codes~\cite{Bravyi2024Memory} are constructed from polynomials over the group algebra $\F_2[x,y]/(x^\ell - 1, y^m - 1)$. Polynomials $A(x,y)$ and $B(x,y)$ define circulant block matrices, yielding:
\begin{equation}
\Hmat_x = [\mathbf{A} \mid \mathbf{B}], \quad \Hmat_z = [\mathbf{B}^\top \mid \mathbf{A}^\top]
\label{eq:bb}
\end{equation}
with $n = 2\ell m$ physical qubits. The CSS condition~\eqref{eq:css} is automatically satisfied. This family achieves an error threshold of 0.7\%---on par with the surface code---while encoding $k = 12$ logical qubits at rates more than an order of magnitude better than the surface code at comparable distances~\cite{Bravyi2024Memory}. Each stabilizer has bounded weight~6, and qubit connectivity can be routed on two physical layers, making BB codes attractive for near-term hardware.

\subsection{Belief propagation on Tanner graphs}\label{sec:bp}

The Tanner graph of a parity-check matrix~$\Hmat$ is a bipartite graph $\mathcal{G} = (\mathcal{V} \cup \mathcal{C}, \mathcal{E})$ connecting variable nodes~$\mathcal{V}$ (qubits) to check nodes~$\mathcal{C}$ (stabilizers). In min-sum BP, decoding proceeds iteratively. At iteration~$t$, the variable-to-check message from node~$i$ to check~$j$ is:
\begin{equation}
\mu_{v \to c}^{(t)}(i,j) = \llr_i + \!\!\sum_{j' \in \mathcal{N}(i)\setminus j}\!\! \mu_{c \to v}^{(t-1)}(j',i)
\label{eq:vtc}
\end{equation}
where $\llr_i = \ln\bigl((1-p_i)/p_i\bigr)$ is the channel log-likelihood ratio. The check-to-variable message is:
\begin{equation}
\mu_{c \to v}^{(t)}(j,i) = \!\!\prod_{i' \in \mathcal{N}(j)\setminus i}\!\! \mathrm{sgn}\bigl(\mu_{v \to c}^{(t)}\bigr) \cdot \!\!\min_{i' \in \mathcal{N}(j)\setminus i}\!\! \bigl|\mu_{v \to c}^{(t)}\bigr|
\label{eq:ctv}
\end{equation}
After $T$ iterations, hard decisions are obtained as $\hat{e}_i = \mathbbm{1}[\llr_i^{\text{total}} < 0]$. BP is provably optimal on tree-structured factor graphs and empirically near-optimal when the Tanner graph has large girth~\cite{Richardson2001Capacity}.

\subsection{The short-cycle problem in QLDPC codes}\label{sec:shortcycle}

The CSS constraint~\eqref{eq:css} forces the existence of short 4- and 6-cycles in QLDPC Tanner graphs. These short cycles cause BP messages to become statistically dependent within a few iterations, violating the conditional independence assumption underlying BP's correctness. The resulting failure mode is oscillatory non-convergence: BP alternates between candidate solutions without settling, producing unreliable hard decisions. Roffe et al.~\cite{Roffe2020Landscape} report convergence rates below 50\% for moderate error rates on the $\code{72}{12}{6}$ code. Post-processing methods including ordered statistics decoding (OSD)~\cite{Fossorier1995OSD}, localized statistics decoding (LSD)~\cite{Panteleev2021QLDPC}, and union-find~\cite{Delfosse2021UF} can recover many unconverged cases but introduce worst-case complexity of $O(n^3)$ or higher.

This systematic BP failure---driven by code structure rather than information-theoretic limits---is the primary motivation for the GNN-based correction approach developed in this work.

\subsection{Noise models}\label{sec:noise}

All experiments employ $Z$-biased depolarizing noise with bias ratio $\eta = 20$:
\begin{equation}
p_z = \frac{p \cdot \eta}{\eta + 1}, \quad p_x = \frac{p}{\eta + 1}
\label{eq:bias}
\end{equation}
This bias ratio is motivated by measurements on superconducting transmon qubits, where $T_1$ relaxation dominates dephasing~\cite{Tuckett2018Bias,Darmawan2021XZZX}. Channel LLRs are set as $\llr_z = \ln\bigl((1-p)/p_z\bigr)$ and $\llr_x = \ln\bigl((1-p)/p_x\bigr)$. The $X$ and $Z$ components are decoded independently, as required by the CSS structure.

To evaluate robustness to noise non-stationarity, we modulate the base error rate $p_\mathrm{base}$ per-shot using three drift models, each grounded in a distinct physical mechanism:

\begin{itemize}[leftmargin=*,topsep=2pt,itemsep=2pt]
\item \textbf{Sinusoidal:} $p(t) = p_\mathrm{base} + A\sin(2\pi t / T)$ with $A = 0.015$, $T = 200$, modeling periodic recalibration cycles.

\item \textbf{Ornstein--Uhlenbeck (OU):} $\mathrm{d}p = \theta(p_\mathrm{base} - p)\,\mathrm{d}t + \sigma\,\mathrm{d}W$ with $\theta = 0.1$, $\sigma = 0.005$, a mean-reverting stochastic process modeling slow environmental drift.

\item \textbf{Random telegraph noise (RTN):} $p(t) \in \{p_\mathrm{base} \pm \delta\}$ with $\delta = 0.01$, switching probability $q = 0.005$ per shot, modeling discrete two-level system (TLS) fluctuator events~\cite{Klimov2018,Burnett2019Decoherence}.
\end{itemize}

Critically, the decoder receives \emph{no direct access} to instantaneous noise parameters; the FiLM module must infer the operating point from the syndrome alone.


% ============================================================
%  III. METHODS
% ============================================================
\section{Methods}\label{sec:methods}

We describe each component of the proposed framework, illustrated in Fig.~\ref{fig:architecture}. The system processes a syndrome through four stages: GNN initial correction, neural BP Stage~1, GNN mid-correction using intermediate marginals, and neural BP Stage~2.

\subsection{Tanner graph representation}\label{sec:tanner}

We construct a unified Tanner graph containing edges derived from both $\Hmat_x$ and $\Hmat_z$, distinguished by a binary edge type indicator ($\tau = 0$ for $\Hmat_x$ edges, $\tau = 1$ for $\Hmat_z$ edges). This enables the GNN to learn distinct message functions for the two stabilizer types while sharing a common graph topology. Each node receives a 4-dimensional feature vector encoding its role and local syndrome information. For variable node~$v_i$:
\begin{equation}
\bvec{x}_i = [\llr_i,\; f_i,\; 1,\; 0]
\label{eq:varfeat}
\end{equation}
where $\llr_i$ is the channel LLR and $f_i$ is the fraction of neighboring check nodes with triggered syndromes. For check node~$c_j$:
\begin{equation}
\bvec{x}_j = [0,\; s_j,\; 0,\; 1]
\label{eq:chkfeat}
\end{equation}
where $s_j \in \{0,1\}$ is the syndrome bit. The final two features serve as one-hot node-type indicators.

\subsection{TannerGNN architecture}\label{sec:tannergnn}

\textsc{TannerGNN} is a message-passing neural network designed for the bipartite, edge-typed structure of CSS Tanner graphs. It consists of input projection, $L$ message-passing layers, and readout.

\paragraph{Input projection.} A linear layer followed by LayerNorm~\cite{Ba2016LayerNorm} maps the heterogeneous 4D features to a shared hidden dimension~$d_h$:
\begin{equation}
\bvec{h}_i^{(0)} = \mathrm{LayerNorm}\bigl(\Wmat_\mathrm{proj}\,\bvec{x}_i + \bvec{b}_\mathrm{proj}\bigr)
\label{eq:proj}
\end{equation}
LayerNorm is essential because node features span disparate scales: LLR values can exceed $|\llr| > 20$, while indicator features are binary.

\paragraph{Message passing.} At layer~$\ell$, each node~$i$ aggregates messages from its neighborhood. Following the principle of edge-type-aware processing~\cite{Ninkovic2024GNN}, we use separate message functions for each edge type:
\begin{align}
\bvec{m}_i^{(\ell)} &= \sum_{j \in \mathcal{N}(i)} \mathrm{MLP}_{e,\tau(i,j)}^{(\ell)}\!\bigl(\bvec{h}_j^{(\ell)}\bigr) \label{eq:msg}\\
\bvec{h}_i^{(\ell+1)} &= \mathrm{MLP}_u^{(\ell)}\!\bigl(\bvec{h}_i^{(\ell)} \concat \bvec{m}_i^{(\ell)}\bigr) \label{eq:update}
\end{align}
where $\mathrm{MLP}_{e,\tau}^{(\ell)}$ denotes the edge-type-specific message network. This allows the network to learn qualitatively different update rules for $X$- and $Z$-stabilizer interactions. Optional enhancements include residual connections ($\bvec{h}_i^{(\ell+1)} \mathrel{+}= \bvec{h}_i^{(\ell)}$), layer normalization, and per-edge attention weights:
\begin{equation}
\alpha_{ij} = \sigma\!\bigl(\Wmat_a [\bvec{h}_i^{(\ell)} \concat \bvec{h}_j^{(\ell)}]\bigr)
\label{eq:attn}
\end{equation}

\paragraph{Readout.} A 2-layer MLP maps each variable node's final hidden state to a scalar LLR correction:
\begin{equation}
\Delta_i = \mathrm{MLP}_\mathrm{read}\!\bigl(\bvec{h}_i^{(L)}\bigr), \quad \llr_i' = \llr_i + \Delta_i
\label{eq:readout}
\end{equation}
The additive correction preserves channel information while allowing the GNN to compensate for structural biases induced by short cycles.

\begin{figure*}[!t]
\centering
\fbox{\parbox{0.95\textwidth}{\centering\vspace{4pt}
\pending{FIGURE 1: Insert architecture\_diagram.pdf --- Full pipeline showing FiLM conditioning path (top) and interleaved decoding pipeline (bottom). Blue = GNN, Green = Neural BP, dashed red = FiLM.}
\vspace{4pt}}}
\caption{Overview of the interleaved GNN-BP decoder with FiLM conditioning. The upper path generates noise-adaptive modulation parameters from a syndrome-derived noise estimate. The lower path decodes through four stages: GNN initial correction, neural BP Stage~1, GNN mid-correction on intermediate marginals, and neural BP Stage~2.}
\label{fig:architecture}
\end{figure*}

\subsection{FiLM conditioning for noise adaptivity}\label{sec:film}

Standard BP assumes fixed, known channel parameters. When the physical error rate drifts---as occurs on real hardware due to TLS fluctuators and calibration decay~\cite{Klimov2018,Burnett2019Decoherence}---fixed-parameter decoders face a mismatch between assumed and actual noise. Feature-wise Linear Modulation (FiLM)~\cite{Perez2018FiLM}, originally proposed for visual reasoning, addresses this by conditioning network activations on an external context signal.

A FiLM generator network maps a scalar noise estimate~$\phat$ to per-layer affine modulation parameters:
\begin{equation}
\bigl(\gamma^{(\ell)},\, \beta^{(\ell)}\bigr) = \mathrm{FiLMGen}(\phat), \quad \ell = 1,\ldots,L
\label{eq:filmgen}
\end{equation}
where $\gamma^{(\ell)}, \beta^{(\ell)} \in \R^{d_h}$. Each GNN layer's output is modulated as:
\begin{equation}
\tilde{\bvec{h}}^{(\ell)} = \gamma^{(\ell)} \odot \bvec{h}^{(\ell)} + \beta^{(\ell)}
\label{eq:film}
\end{equation}
The noise estimate~$\phat$ is derived from the syndrome weight---the fraction of triggered stabilizer measurements---which correlates with the instantaneous error rate and requires no external calibration. During training, each mini-batch samples $p$ uniformly from $[0.01, 0.10]$, exposing the FiLM generator to a broad operating range. At inference, the same model adapts to the prevailing noise rate. To our knowledge, this constitutes the first application of FiLM conditioning to QLDPC decoding; prior work~\cite{Stein2026FiLM} has applied FiLM to repetition codes.

\subsection{Differentiable neural belief propagation}\label{sec:neuralbp}

Following Nachmani et al.~\cite{Nachmani2016Learning,Nachmani2018Deep}, we generalize standard BP by introducing three learnable scalar weights per iteration~$t$:
\begin{align}
\llr_i^{(t,\mathrm{sc})} &= w_\mathrm{ch}^{(t)} \cdot \llr_i \label{eq:wch}\\
\mu_{v \to c}^{(t,\mathrm{sc})} &= w_\mathrm{vtc}^{(t)} \cdot \mu_{v \to c}^{(t)} \label{eq:wvtc}\\
\mu_{c \to v}^{(t,\mathrm{sc})} &= w_\mathrm{ctv}^{(t)} \cdot \mu_{c \to v}^{(t)} \label{eq:wctv}
\end{align}
Weights are parameterized as $w = \mathrm{softplus}(w_\mathrm{raw})$ to ensure strict positivity, with $w_\mathrm{raw}$ initialized to zero so that $w \approx \ln 2 \approx 0.693$ at initialization, yielding near-standard BP. A learnable damping factor $\alpha \in (0,1)$ via the sigmoid function blends successive check-to-variable messages:
\begin{equation}
\mu_{c \to v}^{(t)} \leftarrow \alpha \cdot \mu_{c \to v}^{(t,\mathrm{new})} + (1-\alpha) \cdot \mu_{c \to v}^{(t-1)}
\label{eq:damp}
\end{equation}
The check-to-variable update employs TopK-2 selection for efficiency. The final marginal probabilities are:
\begin{equation}
p_i = \sigma\!\bigl(-\llr_i^{\mathrm{total}}\bigr) = \frac{1}{1 + \exp(\llr_i^{\mathrm{total}})}
\label{eq:marginals}
\end{equation}
Since~\eqref{eq:marginals} is composed entirely of differentiable operations, gradients flow through all $T$~BP iterations back into the GNN, enabling end-to-end training via backpropagation. This is a key advantage over reinforcement learning: the GNN receives direct per-qubit gradient information rather than a single scalar reward, yielding substantially more efficient training.

\subsection{Interleaved GNN-BP architecture}\label{sec:interleaved}

The central architectural contribution is an interleaved pipeline that provides the GNN with two correction opportunities---before and during BP execution. This is motivated by the observation~\cite{Gong2024GNN} that a single pre-BP correction cannot anticipate all failure modes emerging during iterative decoding.

\begin{algorithm}[t]
\caption{Interleaved GNN-BP Decoding}\label{alg:interleaved}
\KwIn{Syndrome $\synd$, channel LLR $\llr$, noise estimate $\phat$}
\KwOut{Error estimate $\ehat$}
$(\gamma, \beta) \leftarrow \mathrm{FiLMGen}(\phat)$\tcp*{Noise conditioning}
$\Delta^{(1)} \leftarrow \textsc{TannerGNN}_1(\llr, \synd, \gamma, \beta)$\;
$\llr' \leftarrow \llr + \Delta^{(1)}$\tcp*{Initial correction}
$\llr_\mathrm{mid} \leftarrow \mathrm{NeuralBP}(\llr', \Hmat, T_1)$\tcp*{BP Stage 1}
$p_i^\mathrm{mid} \leftarrow \sigma(-\llr_{\mathrm{mid},i}) \;\; \forall\, i$\tcp*{Mid-marginals}
$\Delta^{(2)} \leftarrow \textsc{TannerGNN}_2(\llr_\mathrm{mid}, p^\mathrm{mid}, \synd, \gamma, \beta)$\;
$\llr'' \leftarrow \llr_\mathrm{mid} + \Delta^{(2)}$\tcp*{Mid-correction}
$\llr_\mathrm{final} \leftarrow \mathrm{NeuralBP}(\llr'', \Hmat, T_2)$\tcp*{BP Stage 2}
$\hat{e}_i \leftarrow \mathbbm{1}[\llr_{\mathrm{final},i} < 0] \;\; \forall\, i$\;
\Return{$\ehat$}
\end{algorithm}

The procedure (Algorithm~\ref{alg:interleaved}) proceeds in four stages:

\paragraph{Stage 1: GNN initial correction.} The GNN processes channel LLRs and syndrome to produce additive corrections: $\llr' = \llr + \Delta_\mathrm{GNN}^{(1)}$.

\paragraph{Stage 2: Neural BP, first pass.} $T_1$ iterations of neural BP on $\llr'$ produce intermediate marginals $p_i^{\mathrm{mid}} = \sigma(-\llr_{\mathrm{mid},i})$.

\paragraph{Stage 3: GNN mid-correction.} Qubits with $p_i^{\mathrm{mid}} \approx 0.5$ exhibit maximum posterior uncertainty, indicating BP failure. The GNN receives updated features $\bvec{x}_i^{\mathrm{mid}} = [\llr_{\mathrm{mid},i},\; p_i^{\mathrm{mid}},\; 1,\; 0]$ and produces: $\llr'' = \llr_\mathrm{mid} + \Delta_\mathrm{GNN}^{(2)}$.

\paragraph{Stage 4: Neural BP, second pass.} $T_2$ iterations on $\llr''$ yield final hard decisions. The entire pipeline is differentiable end-to-end.

\subsection{Training procedure}\label{sec:training}

\paragraph{Loss function.} The extreme class imbalance in QEC---typically fewer than 5\% of qubits carry errors---motivates focal loss~\cite{Lin2017Focal}:
\begin{equation}
\mathcal{L}_\mathrm{focal} = -\alpha_t\,(1-\hat{p}_t)^\gamma \log(\hat{p}_t)
\label{eq:focal}
\end{equation}
with $\alpha = 0.25$, $\gamma = 2.0$. An auxiliary syndrome consistency loss provides structural signal:
\begin{equation}
\mathcal{L}_\mathrm{syn} = \bigl\|\sigma(\Hmat \cdot \bvec{p}) - \synd\bigr\|^2
\label{eq:synloss}
\end{equation}
Total loss: $\mathcal{L} = \mathcal{L}_\mathrm{focal} + \lambda_\mathrm{syn} \mathcal{L}_\mathrm{syn}$ with $\lambda_\mathrm{syn} = 0.1$.

\paragraph{Optimization.} AdamW~\cite{Loshchilov2019AdamW} with weight decay $10^{-4}$, learning rate $10^{-4}$, cosine annealing with linear warmup. AMP on GPU. Gradient accumulation over 4~mini-batches yields effective batch size~64.

\paragraph{Online data regeneration.} Training data is regenerated each epoch with $p \sim \mathrm{Uniform}[0.01, 0.10]$, preventing overfitting and covering the FiLM operating range.

\subsection{Baseline decoders}\label{sec:baselines}

We compare nine configurations:
\begin{enumerate}[leftmargin=*,topsep=2pt,itemsep=1pt,label=(\arabic*)]
\item \textbf{Plain BP}: min-sum, damping $\alpha = 0.75$, 100 iterations.
\item \textbf{BP-OSD}~\cite{Fossorier1995OSD,Roffe2020Landscape}: BP + order-0 OSD.
\item \textbf{BP-LSD}~\cite{Panteleev2021QLDPC}: BP + localized statistics decoding.
\item \textbf{BeliefFind}~\cite{Delfosse2021UF}: BP + union-find.
\item \textbf{MWPM}~\cite{Higgott2022PyMatching}: minimum-weight perfect matching (approximate for QLDPC).
\item \textbf{GNN-BP}: \textsc{TannerGNN} correction + plain BP.
\item \textbf{GNN+BP-OSD}: \textsc{TannerGNN} correction + BP-OSD.
\item \textbf{GNN+BP-LSD}: \textsc{TannerGNN} correction + BP-LSD.
\item \textbf{Interleaved GNN-BP}: full pipeline (Algorithm~\ref{alg:interleaved}).
\end{enumerate}


% ============================================================
%  IV. EXPERIMENTAL SETUP
% ============================================================
\section{Experimental setup}\label{sec:setup}

\subsection{Code parameters}

\begin{table}[t]
\centering
\caption{Bivariate bicycle codes used in this study.}
\label{tab:codes}
\begin{tabular}{@{}lcccccc@{}}
\toprule
Code & $n$ & $k$ & $d$ & Rate & $\ell$ & $m$ \\
\midrule
BB-72  & 72  & 12 & 6  & 1/6  & 6  & 6  \\
BB-144 & 144 & 12 & 12 & 1/12 & 12 & 6  \\
BB-288 & 288 & 12 & 18 & 1/24 & 12 & 12 \\
\bottomrule
\end{tabular}
\end{table}

Table~\ref{tab:codes} summarizes the three codes. All share the polynomial structure of~\cite{Bravyi2024Memory} with $k=12$ logical qubits and weight-6 stabilizers. Syndrome extraction and error sampling are simulated using Stim~\cite{Gidney2021Stim}.

\subsection{Hyperparameters}

\begin{table}[t]
\centering
\caption{Architecture and training hyperparameters.}
\label{tab:hparams}
\begin{tabular}{@{}ll@{}}
\toprule
\multicolumn{2}{@{}l}{\textit{GNN Architecture}} \\
\midrule
Hidden dimension $d_h$ & 64 \\
Message-passing layers $L$ & 3 \\
Edge types & 2 ($\Hmat_x$, $\Hmat_z$) \\
Node feature dimension & 4 \\
Dropout & 0.1 \\
Correction mode & Additive \\
\midrule
\multicolumn{2}{@{}l}{\textit{Training}} \\
\midrule
Optimizer & AdamW \\
Learning rate & $10^{-4}$ \\
Weight decay & $10^{-4}$ \\
Schedule & Cosine + linear warmup \\
Effective batch size & 64 \\
Focal loss $(\alpha, \gamma)$ & $(0.25, 2.0)$ \\
Syndrome loss weight & 0.1 \\
\midrule
\multicolumn{2}{@{}l}{\textit{BP}} \\
\midrule
Iterations (train / eval) & 10 / 100 \\
Damping & 0.75 (or learnable) \\
Interleaved split $(T_1, T_2)$ & (5, 5) / (50, 50) \\
\bottomrule
\end{tabular}
\end{table}

\subsection{Evaluation metrics}

\textbf{Logical error rate} (LER): fraction of shots with any incorrect logical observable. Wilson score 95\% confidence intervals.

\textbf{BP convergence rate}: fraction of shots where all parity checks are satisfied.

\textbf{McNemar's paired test}~\cite{McNemar1947}: for binary decoder outcomes,
\begin{equation}
\chi^2 = \frac{(|b - c| - 1)^2}{b + c}
\label{eq:mcnemar}
\end{equation}
where $b$ counts shots where decoder~A succeeds and B fails, $c$ the converse. Significance at $p < 0.05$.

\textbf{Wall-clock time}: mean per-shot latency on a single CPU core.


% ============================================================
%  V. RESULTS
% ============================================================
\section{Results}\label{sec:results}

\subsection{Decoder comparison under static noise}

\begin{figure}[t]
\centering
\fbox{\parbox{0.9\columnwidth}{\centering\vspace{12pt}
\pending{FIGURE 2: Bar chart of LER for all 9 decoders on $\code{72}{12}{6}$ at $p=0.02$, $\eta=20$, 5000 shots. Wilson 95\% CI error bars.}
\vspace{12pt}}}
\caption{Logical error rate across nine decoders on $\code{72}{12}{6}$ at $p = 0.02$ under $Z$-biased noise with $\eta=20$.}
\label{fig:bars}
\end{figure}

Figure~\ref{fig:bars} presents the logical error rate achieved by each decoder on the $\code{72}{12}{6}$ code under static noise at $p = 0.02$. \pending{DATA: Insert exact LER values and CIs.}

Three observations are noteworthy. First, \pending{DATA: GNN-BP vs Plain BP --- report LER reduction percentage}. Second, \pending{DATA: GNN+BP-OSD vs BP-OSD --- report whether GNN pre-correction reduces OSD burden}. Third, \pending{DATA: Interleaved vs single-pass GNN-BP --- report additional gain from mid-correction}.

The improvement stems from two complementary mechanisms: the GNN's additive LLR corrections break short-cycle correlations that cause BP oscillation, while the neural BP weights learn iteration-specific scaling that suppresses remaining oscillatory modes.

\subsection{Threshold behavior}

\begin{figure}[t]
\centering
\fbox{\parbox{0.9\columnwidth}{\centering\vspace{12pt}
\pending{FIGURE 3: Semilog plot of LER vs.\ $p \in [0.005, 0.10]$ for selected decoders on $\code{72}{12}{6}$.}
\vspace{12pt}}}
\caption{Logical error rate vs.\ physical error rate.}
\label{fig:threshold}
\end{figure}

Figure~\ref{fig:threshold} shows LER as a function of physical error rate. \pending{DATA: Insert curves and describe threshold behavior.}

Two regimes are apparent. At low error rates ($p < 0.01$), \pending{DATA: all decoders perform comparably}. As $p$ increases beyond $\approx 0.015$, \pending{DATA: plain BP plateaus near LER~$\approx 0.5$ due to convergence failure, while GNN-corrected decoders maintain a steep descent, extending BP's usable operating range}.

\subsection{BP convergence analysis}

\begin{figure}[t]
\centering
\fbox{\parbox{0.9\columnwidth}{\centering\vspace{12pt}
\pending{FIGURE 4: BP convergence rate vs.\ $p$ for Plain BP (dashed) and GNN-BP (solid), all three codes.}
\vspace{12pt}}}
\caption{BP convergence rate vs.\ physical error rate.}
\label{fig:convergence}
\end{figure}

Figure~\ref{fig:convergence} plots BP convergence rate as a function of $p$. \pending{DATA: Insert rates for all three codes.}

\pending{DATA: Plain BP convergence drops below 50\% at $p \approx 0.02$; GNN-corrected BP maintains convergence above 90\%.} This confirms that BP failure on QLDPC codes is a correctable artifact of poor initial priors rather than an intrinsic algorithmic limitation.

\subsection{Robustness under temporal drift}

\begin{figure}[t]
\centering
\fbox{\parbox{0.9\columnwidth}{\centering\vspace{12pt}
\pending{FIGURE 5: Grouped bars for LER under static vs.\ 3 drift models for Plain BP, GNN-BP (no FiLM), GNN-BP (with FiLM) at $p_\mathrm{base}=0.02$.}
\vspace{12pt}}}
\caption{Decoder performance under static and drifting noise.}
\label{fig:drift}
\end{figure}

Figure~\ref{fig:drift} evaluates robustness to temporal drift. \pending{DATA: Insert LER values under each drift model.}

\pending{DATA: The critical comparison is GNN-BP with vs.\ without FiLM. FiLM-conditioned decoders are expected to maintain near-static performance, demonstrating that syndrome-derived noise estimation provides sufficient signal for real-time adaptation.}

\subsection{Multi-code scaling}

\begin{figure}[t]
\centering
\fbox{\parbox{0.9\columnwidth}{\centering\vspace{12pt}
\pending{FIGURE 6: LER vs.\ $p$ for all three code sizes.}
\vspace{12pt}}}
\caption{Scaling behavior across three BB code sizes.}
\label{fig:scaling}
\end{figure}

Figure~\ref{fig:scaling} examines how GNN correction scales with code size. \pending{DATA: Insert curves and report whether relative improvement is maintained.}

\subsection{Component ablation}

\begin{table}[t]
\centering
\caption{Component ablation on $\code{72}{12}{6}$ at $p = 0.02$.}
\label{tab:ablation}
\begin{tabular}{@{}lccc@{}}
\toprule
Configuration & LER & $\Delta$LER & Conv. \\
\midrule
Plain BP & \pending{---} & --- & \pending{---} \\
\; + Neural BP weights & \pending{---} & \pending{---} & \pending{---} \\
\; + GNN correction & \pending{---} & \pending{---} & \pending{---} \\
\; + FiLM conditioning & \pending{---} & \pending{---} & \pending{---} \\
\; + Attention & \pending{---} & \pending{---} & \pending{---} \\
\; + Interleaved & \pending{---} & \pending{---} & \pending{---} \\
\bottomrule
\end{tabular}
\end{table}

Table~\ref{tab:ablation} isolates each component's marginal contribution. \pending{DATA: Fill and describe which components provide largest gains.}

\subsection{Transfer learning}

\begin{table}[t]
\centering
\caption{Transfer learning from $\code{72}{12}{6}$ to larger codes.}
\label{tab:transfer}
\begin{tabular}{@{}llcc@{}}
\toprule
Target & Strategy & LER & vs.\ Scratch \\
\midrule
\multirow{2}{*}{$\code{144}{12}{12}$} & From scratch & \pending{---} & --- \\
& Frozen transfer & \pending{---} & \pending{---} \\
\midrule
\multirow{2}{*}{$\code{288}{12}{18}$} & From scratch & \pending{---} & --- \\
& Frozen transfer & \pending{---} & \pending{---} \\
\bottomrule
\end{tabular}
\end{table}

Table~\ref{tab:transfer} evaluates frozen-backbone transfer. \pending{DATA: Fill table and report whether local features transfer.}

\subsection{Computational cost}

\begin{table}[t]
\centering
\caption{Mean per-shot decoding latency (single CPU core).}
\label{tab:timing}
\begin{tabular}{@{}lcc@{}}
\toprule
Decoder & Time (ms) & Relative \\
\midrule
Plain BP & \pending{---} & $1.0\times$ \\
GNN-BP & \pending{---} & \pending{---} \\
BP-OSD & \pending{---} & \pending{---} \\
Interleaved GNN-BP & \pending{---} & \pending{---} \\
MWPM & \pending{---} & \pending{---} \\
\bottomrule
\end{tabular}
\end{table}

Table~\ref{tab:timing} reports decoding latency. \pending{DATA: Fill table.}

\subsection{Statistical significance}

\begin{table}[t]
\centering
\caption{McNemar's test vs.\ plain BP ($\code{72}{12}{6}$, $p=0.02$).}
\label{tab:mcnemar}
\begin{tabular}{@{}lccc@{}}
\toprule
Decoder & $\chi^2$ & $p$-value & Sig. \\
\midrule
GNN-BP & \pending{---} & \pending{---} & \pending{---} \\
BP-OSD & \pending{---} & \pending{---} & \pending{---} \\
GNN+BP-OSD & \pending{---} & \pending{---} & \pending{---} \\
Interleaved & \pending{---} & \pending{---} & \pending{---} \\
\bottomrule
\end{tabular}
\end{table}

Table~\ref{tab:mcnemar} reports paired significance tests. \pending{DATA: Fill table.}


% ============================================================
%  VI. DISCUSSION
% ============================================================
\section{Discussion}\label{sec:discussion}

\paragraph{Principal finding.}
The results demonstrate that augmenting BP with learned LLR corrections is more effective than either replacing BP entirely or relying on expensive post-processing. The GNN addresses the specific failure mode---short-cycle-induced message correlation---while preserving BP's convergence guarantees and $O(n)$ per-iteration complexity.

\paragraph{Comparison with Astra.}
Maan and Paler~\cite{Maan2025MLBP} replace BP entirely with a learned GNN, achieving higher thresholds than BP-OSD on BB codes up to distance~18 under depolarizing phenomenological noise. Our approach differs in three respects: (i)~we augment rather than replace BP, preserving compatibility with existing post-processing; (ii)~we evaluate under $Z$-biased noise ($\eta = 20$), more representative of superconducting hardware; and (iii)~we address temporal drift. The two approaches are complementary: GNN correction could be combined with Astra's architecture as an alternative to standard BP.

\paragraph{Comparison with differentiable iterative decoders.}
Gong et al.~\cite{Gong2024GNN} interleave BP with GNN layers for quantum bicycle codes, demonstrating that mid-decoding intervention lowers the error floor. Our architecture extends this with FiLM conditioning, neural BP weights, and evaluation on the larger BB family under biased and drifting noise. The independent development of interleaving strategies in~\cite{Gong2024GNN} and the present work suggests mid-decoding correction is a broadly useful principle.

\paragraph{Comparison with FiLM-conditioned decoders.}
Stein et al.~\cite{Stein2026FiLM} apply FiLM to repetition-code decoders. We extend this to QLDPC codes, where irregular Tanner graphs and the CSS constraint create a substantially harder decoding landscape. Our syndrome-weight noise estimation is simpler than calibration-derived estimates but generalizes across drift models.

\paragraph{Why not reinforcement learning?}
The RL formulation (agent adjusts BP priors, receives binary reward) suffers from severe credit assignment: a single success/failure signal provides no information about \emph{which} of 72+ qubit priors were wrong. End-to-end differentiable training through BP provides per-qubit gradients via focal loss, yielding orders-of-magnitude better sample efficiency. FiLM conditioning handles noise adaptation that RL would otherwise require online exploration to achieve.

\paragraph{Limitations.}
(i)~All experiments use Stim simulation; hardware validation is needed. (ii)~MWPM is approximate for QLDPC. (iii)~Code-capacity noise omits measurement errors and leakage. (iv)~GNN latency may challenge real-time $\mu$s-scale decoding. (v)~Syndrome-weight noise estimation may degrade at very high error rates.

\paragraph{Design philosophy.}
BP is a powerful algorithm whose failure on QLDPC codes is caused by a specific, correctable pathology rather than a fundamental limitation. The GNN corrects channel priors to account for code structure, noise bias, and drift, while BP handles structured inference. This modular design enables independent component updates and integration with existing decoder infrastructure.


% ============================================================
%  VII. CONCLUSION
% ============================================================
\section{Conclusion}\label{sec:conclusion}

We have presented a hybrid decoding framework for quantum LDPC codes that integrates graph neural networks with differentiable belief propagation in a novel interleaved architecture. The four components---\textsc{TannerGNN} with edge-type-aware message passing, FiLM conditioning for noise adaptivity, learnable BP message weights, and mid-decoding GNN intervention---together \pending{DATA: achieve X\% LER reduction relative to standard BP on $\code{72}{12}{6}$, improving convergence from below 50\% to above 90\%}. The framework maintains performance under temporal noise drift and transfers to larger codes within the BB family.

Promising directions for future work include validation on hardware syndrome data from IBM and Google processors, scaling to codes with thousands of qubits, and FPGA-based quantized inference for real-time deployment. The augment-not-replace principle explored here may generalize to other quantum code families with challenging Tanner graph structures, including lifted-product and quantum expander codes.


% ============================================================
%  ACKNOWLEDGMENTS
% ============================================================
\section*{Acknowledgments}

The authors thank Dr.~Daryn Stover for mentorship and the staff of the STAR Lab at the University of Arizona for providing guidance and computational resources.


% ============================================================
%  DATA AVAILABILITY
% ============================================================
\section*{Data and code availability}

Source code for \textsc{TannerGNN}, the training pipeline, all decoder implementations, and reproduction scripts are available at \pending{URL: Insert GitHub repository}. Simulated datasets were generated using Stim~\cite{Gidney2021Stim}.


% ============================================================
%  REFERENCES
% ============================================================
\bibliographystyle{quantum}
\bibliography{references}

\end{document}
